\chapter{System Requirements Specification}
\label{chap:srs}

\section{Functional Requirements}
\label{sec:functional_requirements}

\subsection{FR-1: User Authentication and Account Management}

\begin{itemize}
  \item The system shall allow users to register and log in using email and password.
  \item The system shall support Google OAuth~2.0 via PKCE flow for mobile single-sign-on.
  \item The system shall support secure password reset via email.
  \item The system shall enforce email confirmation before granting full account access.
  \item The system shall collect a unique username during first-time onboarding.
  \item The system shall allow users to permanently delete their accounts, with cascading deletion of all associated data.
\end{itemize}

\subsection{FR-2: Video Upload and Processing}

\begin{itemize}
  \item The system shall accept video uploads from the device camera roll or live recording.
  \item The system shall enforce a maximum upload size of 200\,MB.
  \item The system shall detect and reject videos that are too short (fewer than 30 frames) to support reliable analysis.
  \item The system shall allow users to specify or override the automatically detected shooting side (left/right).
  \item The system shall queue submitted videos and process them asynchronously, returning a job identifier immediately.
  \item The system shall cap concurrent analysis jobs at eight globally, returning HTTP~429 to excess requests.
\end{itemize}

\subsection{FR-3: Biomechanical Analysis}

\begin{itemize}
  \item The system shall extract per-frame 33-landmark body pose keypoints using MediaPipe Pose.
  \item The system shall smooth noisy pose signals using Savitzky--Golay filtering.
  \item The system shall detect the shot window (crouch, release, and landing events) using a multi-signal state machine.
  \item The system shall compute elbow extension at release, knee flexion at crouch, and wrist follow-through angles.
  \item The system shall optionally compute ball release angle via YOLOv8n trajectory fitting.
  \item The system shall assess each metric against research-validated normative ranges, issuing verdicts of \textit{Good}, \textit{Needs Work}, or \textit{Low Confidence}.
  \item The system shall produce an overall shot quality score on a 0--100 scale.
  \item The system shall generate an annotated overlay video and raw data artefacts (pose CSV, angles CSV, JSON report).
\end{itemize}

\subsection{FR-4: AI Coaching Feedback}

\begin{itemize}
  \item The system shall generate natural-language shot feedback using Gemini 2.5 Flash.
  \item Feedback shall include: an overall explanation, metric-specific verdicts, listed strengths, improvement cues, and concise coaching bullet points.
  \item The system shall provide a holistic AI score (0--100) alongside the rule-based metric score.
  \item The system shall fall back to the rule-based score if Gemini fails or exceeds a 10-second timeout.
  \item The system shall expose a conversational coaching chat interface using streaming server-sent events.
  \item The chat interface shall inject the user's recent analysis history, profile metadata, and weak areas as context for the LLM.
\end{itemize}

\subsection{FR-5: Drill Recommendations}

\begin{itemize}
  \item The system shall recommend drills personalised to the user's identified mechanical weaknesses.
  \item Recommendations shall use FAISS nearest-neighbour similarity search over a curated drill pool.
  \item The system shall employ a LinUCB contextual bandit to rank drill clusters based on the user's mechanical profile.
  \item Each recommendation shall include a drill name, rationale, step-by-step instructions, and estimated duration.
\end{itemize}

\subsection{FR-6: History and Progress Tracking}

\begin{itemize}
  \item The system shall persist completed analyses to the user's cloud account.
  \item The system shall maintain a daily analysis streak counter.
  \item The system shall expose an analysis history endpoint with pagination and filtering.
  \item The system shall display aggregated progress metrics (weekly average score, total analyses, longest streak).
\end{itemize}

\subsection{FR-7: Profile Management}

\begin{itemize}
  \item The system shall maintain user profile data including skill level, position, dominant hand, age, height, and coaching style preference.
  \item Profile data shall be used to personalise normative thresholds, AI feedback tone, and drill difficulty.
\end{itemize}

\section{Non-Functional Requirements}
\label{sec:nfr}

\subsection{NFR-1: Performance}
\begin{itemize}
  \item The video analysis pipeline shall complete within 40 seconds for a 5-second clip at 30 FPS under normal server load.
  \item API responses for non-analysis endpoints shall complete within 500\,ms under normal load.
  \item The frontend shall display analysis results within 3 seconds of job completion notification.
\end{itemize}

\subsection{NFR-2: Scalability}
\begin{itemize}
  \item The backend shall support a minimum of eight concurrent analysis jobs without degradation.
  \item The job concurrency limit shall be configurable via environment variable without code changes.
  \item The architecture shall support horizontal scaling by decoupling the job queue from the HTTP process.
\end{itemize}

\subsection{NFR-3: Usability}
\begin{itemize}
  \item The application shall be operable by a non-technical user without a user manual.
  \item Visual feedback shall be provided for all asynchronous operations lasting more than one second.
  \item Error states shall display user-friendly messages rather than raw exception text.
\end{itemize}

\subsection{NFR-4: Reliability}
\begin{itemize}
  \item Failed Gemini enrichment shall not prevent analysis results from being returned to the client.
  \item Failed result persistence shall not crash the application; retry logic shall be applied.
  \item The system shall degrade gracefully when Supabase is unavailable, continuing to serve cached results where possible.
\end{itemize}

\subsection{NFR-5: Security}
\begin{itemize}
  \item All user-scoped data access shall be protected by Supabase Row-Level Security policies.
  \item API keys shall be stored exclusively in environment variables, never in source code.
  \item Authentication tokens shall be transmitted in request headers, not URL query parameters.
  \item Per-IP rate limiting shall be applied to all upload endpoints.
\end{itemize}

\subsection{NFR-6: Portability}
\begin{itemize}
  \item The mobile application shall run on iOS~16+ and Android~12+ without platform-specific code paths.
  \item The backend shall run on Linux, macOS, and Windows (Python~3.11+) without modification.
\end{itemize}

\subsection{NFR-7: Maintainability}
\begin{itemize}
  \item All API contracts shall be expressed in Pydantic models (backend) and TypeScript interfaces (frontend).
  \item All configuration parameters shall be centralised in \lstinline|mvp_config.yaml| and environment variable files.
  \item All production log messages shall use structured JSON logging via \lstinline|logging.getLogger(__name__)|.
\end{itemize}

\section{User Requirements}
\label{sec:user_requirements}

\begin{itemize}
  \item As a player, I want to record or upload a video of my shot and receive a score so I can gauge my current form.
  \item As a player, I want to see which specific joints are out of alignment and understand why, so I can correct them.
  \item As a player, I want to ask the AI coach follow-up questions about my analysis in natural language.
  \item As a player, I want to be recommended drills that specifically target my weakest areas.
  \item As a player, I want to track my score over time and see whether my mechanics are improving.
  \item As a coach, I want to review individual player analyses and see metric breakdowns so I can provide targeted guidance.
\end{itemize}

\section{System Constraints}
\label{sec:constraints}

\begin{itemize}
  \item Video analysis is CPU-bound and requires a server with multiple cores to meet performance targets; GPU inference is not currently required.
  \item MediaPipe Pose requires that the full body be visible in the frame; heavily occluded or partially visible clips will produce low-confidence results.
  \item Ball trajectory analysis via YOLOv8n is disabled by default. Enabling it requires setting \lstinline|SHOOTRZ_ENABLE_BALL=1|, due to the additional per-analysis CPU cost.
  \item The Gemini 2.5 Flash API has a rate limit of approximately 200 requests per minute at the free/standard tier.
  \item Supabase Storage is used for persisted video artefacts; the free tier imposes a 1\,GB storage limit.
\end{itemize}

\section{Assumptions}
\label{sec:assumptions}

\begin{itemize}
  \item Users record videos from a fixed camera position approximately 3--5 metres from the player, with the full body visible from the side or front.
  \item Videos are in MP4 or a compatible container format; HEVC (the iOS default) may require transcoding before submission.
  \item Users have a stable internet connection during video upload and polling.
  \item The server has at least two available CPU cores for pipeline processing.
\end{itemize}

\section{Use Cases}
\label{sec:use_cases}

\subsection{UC-1: Analyse Shot}

\begin{description}
  \item[Actor] Authenticated user
  \item[Trigger] User selects a video from their camera roll
  \item[Precondition] User is logged in; video is at least 30 frames; fewer than 8 jobs are in flight
  \item[Flow] Upload $\rightarrow$ Queue $\rightarrow$ Pipeline $\rightarrow$ Gemini enrichment $\rightarrow$ Poll $\rightarrow$ Display $\rightarrow$ Persist
  \item[Postcondition] Analysis is saved to user's account; streak counter is updated
\end{description}

\subsection{UC-2: Chat with Coach}

\begin{description}
  \item[Actor] Authenticated user
  \item[Trigger] User types a question in the Chat screen
  \item[Precondition] User has at least one completed analysis in their history
  \item[Flow] Submit message $\rightarrow$ Backend injects analysis context $\rightarrow$ Gemini streams response $\rightarrow$ UI renders incrementally
  \item[Postcondition] Message is persisted to the \lstinline|chat_history| table
\end{description}

\subsection{UC-3: Get Drill Recommendations}

\begin{description}
  \item[Actor] Authenticated user
  \item[Trigger] User taps the Drills screen or ``Get Drills'' from an analysis result
  \item[Flow] Send user metrics vector $\rightarrow$ FAISS search $\rightarrow$ LinUCB arm selection $\rightarrow$ Filter and rank $\rightarrow$ Return drill
  \item[Postcondition] Drill is displayed with rationale and step-by-step instructions
\end{description}

\subsection{UC-4: View Progress}

\begin{description}
  \item[Actor] Authenticated user
  \item[Trigger] User opens the Progress screen
  \item[Flow] Fetch history from Supabase $\rightarrow$ Compute weekly averages, streak, trend $\rightarrow$ Render charts
  \item[Postcondition] Historical performance trends are displayed
\end{description}

\section{User Flow}
\label{sec:user_flow}

Figure~\ref{fig:user_flow} illustrates the high-level navigation and analysis flow for both new and returning users.

\begin{figure}[H]
  \centering
  \figurebox[0.88\textwidth]{images/user_flow.png}
  \caption{SHOOTRZ application user flow: onboarding, authentication, and analysis cycle.}
  \label{fig:user_flow}
\end{figure}

On launching the application, a new user is directed through onboarding slides before completing email or Google authentication and selecting a username. A returning user is directed immediately to the Home Dashboard. From the dashboard, the user can navigate to the Analyse, Chat, Drills, or Progress screens. The analysis flow involves selecting a video, submitting it for processing, polling for the result, and optionally saving the analysis to their account. The full cycle---analyse, review feedback, complete recommended drills, re-analyse---constitutes the core practice loop.
